% ============================================================
% Lezione 01: Lezione Test
% Generato: 2026-02-28 00:00
% Fonti: Light Nature.pptx
% ============================================================

\section{Lezione 1: Lezione Test}

\label{sec:lezione01}


\subsection{Light:its nature and its properties}
\label{subsec:slide01-1}

Light:its nature and its properties

Lecturer: Mauro Mosca
(www.dieet.unipa.it/tfl)
last update: 29/09/2021

University of Palermo – Master in Electronic Engineering


\subsection{Waves}
\label{subsec:slide01-2}

\begin{figure}[H]
\centering
\includegraphics[width=0.8\textwidth]{slide002_310598c7.wmf}
\end{figure}

Longitudinal waves

\begin{figure}[H]
\centering
\includegraphics[width=0.8\textwidth]{slide002_f9ce4266.gif}
\end{figure}

Transversal waves

\begin{figure}[H]
\centering
\includegraphics[width=0.8\textwidth]{slide002_20e596a5.gif}
\end{figure}


\subsection{Frequency and wavelength}
\label{subsec:slide01-3}

\begin{figure}[H]
\centering
\includegraphics[width=0.8\textwidth]{slide003_2b48e4b2.png}
\end{figure}

Not all waves travel…

…at a fixed point in space

\begin{figure}[H]
\centering
\includegraphics[width=0.8\textwidth]{slide003_3057ad1b.png}
\end{figure}

\begin{figure}[H]
\centering
\includegraphics[width=0.8\textwidth]{slide003_3fdced3d.gif}
\end{figure}

standing waves


\subsection{Frequency and wavelength}
\label{subsec:slide01-4}

l

l/T = lf = v


\subsection{Light}
\label{subsec:slide01-5}

particles

or

Waves

?

+

reflection, refraction

shadows

Electromagnetic and quantum theory are in perfect agreement if we assume that the amplitude of a field (electric or magnetic) at a certain point and at a certain instant is proportional to the probability (wave function) of finding a photon

\begin{figure}[H]
\centering
\includegraphics[width=0.8\textwidth]{slide005_87b51c96.gif}
\end{figure}

\begin{figure}[H]
\centering
\includegraphics[width=0.8\textwidth]{slide005_77763c66.jpg}
\end{figure}

\begin{figure}[H]
\centering
\includegraphics[width=0.8\textwidth]{slide005_8863b783.gif}
\end{figure}

\begin{figure}[H]
\centering
\includegraphics[width=0.8\textwidth]{slide005_d2480124.gif}
\end{figure}


\subsection{Light properties}
\label{subsec:slide01-6}

Reflection

qr

qi

\begin{figure}[H]
\centering
\includegraphics[width=0.8\textwidth]{slide006_e85e96fc.wmf}
\end{figure}

qi = qr

n1

n2

qt

n1sinqi = n2sinqt

Refraction


\subsection{Refractive indices}
\label{subsec:slide01-7}

\begin{figure}[H]
\centering

\includegraphics[width=0.8\textwidth]{slide007_24462c13.jpg}
\end{figure}


\subsection{Why refraction?}
\label{subsec:slide01-8}

\begin{figure}[H]
\centering
\includegraphics[width=0.8\textwidth]{slide008_0a658c35.jpg}
\end{figure}

n1 = 1

\begin{figure}[H]
\centering

\includegraphics[width=0.8\textwidth]{slide008_61aac317.png}
\end{figure}

?

n2 = 1,3


\subsection{Interference: some pretty simple math...}
\label{subsec:slide01-9}

\begin{figure}[H]
\centering
\includegraphics[width=0.8\textwidth]{slide009_250ad787.png}
\end{figure}

\begin{figure}[H]
\centering
\includegraphics[width=0.8\textwidth]{slide009_944f3aaa.png}
\end{figure}

asin = n

asin = (n+1/2)


\subsection{Slide 10}
\label{subsec:slide01-10}

\begin{figure}[H]
\centering
\includegraphics[width=0.8\textwidth]{slide010_a2fbb739.png}
\end{figure}

\begin{figure}[H]
\centering
\includegraphics[width=0.8\textwidth]{slide010_41249af0.png}
\end{figure}

\begin{figure}[H]
\centering
\includegraphics[width=0.8\textwidth]{slide010_310598c7.wmf}
\end{figure}

ngsir.netfirms.com

Web project:


\subsection{Photoelectric effect}
\label{subsec:slide01-11}

\begin{equation}
\text{[formula]}
\end{equation}

both velocity and energy of electrons independent of light intensity

both velocity and energy of electrons dependent on color of light

electrons insensitive to some colors


\subsection{Photons}
\label{subsec:slide01-12}

photons

Light = “energy packets”

E = hf = hc/l   (Planck-Einstein relation)

\begin{figure}[H]
\centering
\includegraphics[width=0.8\textwidth]{slide012_eec5fc70.png}
\end{figure}


\subsection{Electromagnetic spectrum}
\label{subsec:slide01-13}

\begin{figure}[H]
\centering
\includegraphics[width=0.8\textwidth]{slide013_93ed602f.png}
\end{figure}

\begin{figure}[H]
\centering
\includegraphics[width=0.8\textwidth]{slide013_f9a25343.png}
\end{figure}

\begin{figure}[H]
\centering
\includegraphics[width=0.8\textwidth]{slide013_ee39dcd4.png}
\end{figure}


\subsection{Dispersion}
\label{subsec:slide01-14}

\begin{figure}[H]
\centering
\includegraphics[width=0.8\textwidth]{slide014_53f1e6b2.jpg}
\end{figure}

\begin{figure}[H]
\centering
\includegraphics[width=0.8\textwidth]{slide014_daf17bbb.gif}
\end{figure}

\begin{figure}[H]
\centering
\includegraphics[width=0.8\textwidth]{slide014_e85e96fc.wmf}
\end{figure}

v  function of l

n = n (l)

\begin{figure}[H]
\centering
\includegraphics[width=0.8\textwidth]{slide014_13ba3ba0.png}
\end{figure}


\subsection{Gratings}
\label{subsec:slide01-15}

\begin{figure}[H]
\centering
\includegraphics[width=0.8\textwidth]{slide015_147d141d.jpg}
\end{figure}

in reflection

in trasmission

\begin{figure}[H]
\centering
\includegraphics[width=0.8\textwidth]{slide015_250ad787.png}
\end{figure}

\begin{figure}[H]
\centering
\includegraphics[width=0.8\textwidth]{slide015_dd37f00f.png}
\end{figure}

\begin{figure}[H]
\centering
\includegraphics[width=0.8\textwidth]{slide015_a6274ef4.wmf}
\end{figure}

asin = n


\subsection{Spectrometer}
\label{subsec:slide01-16}

\begin{figure}[H]
\centering
\includegraphics[width=0.8\textwidth]{slide016_dc0eaa88.png}
\end{figure}

\begin{figure}[H]
\centering
\includegraphics[width=0.8\textwidth]{slide016_c7416999.gif}
\end{figure}

mirrors

\begin{figure}[H]
\centering
\includegraphics[width=0.8\textwidth]{slide016_66ec243c.jpg}
\end{figure}

slit

resolution?…


\subsection{Polarization of light}
\label{subsec:slide01-17}

E

polarized light

H

non-polarized 
light


\subsection{Polarizers}
\label{subsec:slide01-18}

E

\begin{figure}[H]
\centering
\includegraphics[width=0.8\textwidth]{slide018_310598c7.wmf}
\end{figure}


\subsection{Optical activity}
\label{subsec:slide01-19}

ability of some crystals to rotate the wave 
polarization plane

possible in crystals that exhibit 
birefringence


\subsection{Light modulation}
\label{subsec:slide01-20}

\begin{figure}[H]
\centering
\includegraphics[width=0.8\textwidth]{slide020_836c9c46.png}
\end{figure}

\begin{figure}[H]
\centering
\includegraphics[width=0.8\textwidth]{slide020_39cf8b4c.png}
\end{figure}

\begin{equation}
\frac{\cdot}{\cdot}
\end{equation}

\begin{equation}
\frac{\cdot}{\cdot}
\end{equation}

\begin{equation}
\frac{\cdot}{\cdot}
\end{equation}

\begin{equation}
\frac{\cdot}{\cdot}
\end{equation}

\begin{equation}
\frac{\cdot}{\cdot}
\end{equation}


\subsection{Light modulation: amplitude modulation}
\label{subsec:slide01-21}

Light modulation: amplitude modulation

\begin{figure}[H]
\centering
\includegraphics[width=0.8\textwidth]{slide021_d3979c9b.png}
\end{figure}

= f(t)
